% GENERATED FILE. Edit canonical lessons, then run npm run generate:books.
% canonical-source-hash: fnv1a64:f794793a481dfaac
% canonical-lessons: PA-C323-dhidd, PA-C323-haddi, PA-C323-kaleja, PA-C323-pasli, PA-C323-dukhna

\chapter{Inside the Body}
\label{ch:pa-a2-inside-the-body}
\hlchaptermodality{323}

\begin{chapteropening}
\textbf{By the end of this chapter:} \emph{I can say which part inside my body hurts.}

Complete the last lesson of chapter 323: I can say which part inside my body hurts.
\end{chapteropening}

\section[\texorpdfstring{\pa{ਢਿੱਡ} (ḍhiḍḍ)}{ਢਿੱਡ (ḍhiḍḍ)}]{\pa{ਢਿੱਡ} (ḍhiḍḍ) — the belly}

\label{lesson:PA-C323-dhidd}



\begin{quote}
Before the new one: say the Punjabi for a bandage, then the Punjabi for an ointment.
\end{quote}

\subsection*{You'll want to know: \pa{ਢਿੱਡ}}
\textbf{\pa{ਢਿੱਡ}} — \emph{ḍhiḍḍ} — \textquotedblleft{}the belly\textquotedblright{}.

\begin{quote}
\pa{ਮੇਰੇ} \pa{ਢਿੱਡ} \pa{ਵਿੱਚ} \pa{ਦਰਦ} \pa{ਹੈ।}
\end{quote}

\begin{quote}
\emph{mere ḍhiḍḍ vich darad hai.}
\end{quote}

\begin{quote}
I have a pain in my belly.
\end{quote}

A masculine noun.

\begin{cousinweb}
\textbf{\pa{ਢਿੱਡ}} (\emph{ḍhiḍḍ}) is inherited from a reconstructed \emph{*ḍhiḍḍha}, \textquotedblleft{}belly\textquotedblright{}; CDIAL lists Punjabi \emph{ḍhiḍḍh}, \textquotedblleft{}belly\textquotedblright{}, under it (CDIAL 5589).
\end{cousinweb}

\begin{grammarlens}[title={What you've built}]
One more word for this chapter.
\end{grammarlens}

\subsection*{Your turn}
\begin{itemize}
\raggedright
  \item \emph{Say it:} \emph{ḍhiḍḍ}
  \item \emph{Say it:} \emph{ḍhiḍḍ}, once more
  \item \emph{Say it:} \emph{mere ḍhiḍḍ vich darad hai.} — I have a pain in my belly.
  \item \emph{Recall:} say the Punjabi for a bandage, then the Punjabi for an ointment, then say \emph{ḍhiḍḍ} again
\end{itemize}

\begin{tcolorbox}[breakable,colback=teal!4,colframe=teal!35!black,title={Before you move on}]
What does \pa{ਢਿੱਡ} mean? (\textquotedblleft{}The belly\textquotedblright{}.) Say it once more.
\end{tcolorbox}
\section[\texorpdfstring{\pa{ਹੱਡੀ} (haḍḍī)}{ਹੱਡੀ (haḍḍī)}]{\pa{ਹੱਡੀ} (haḍḍī) — a bone}

\label{lesson:PA-C323-haddi}



\begin{quote}
Before the new one: say the Punjabi for an ointment, then the Punjabi for the belly.
\end{quote}

\subsection*{You'll want to know: \pa{ਹੱਡੀ}}
\textbf{\pa{ਹੱਡੀ}} — \emph{haḍḍī} — \textquotedblleft{}a bone\textquotedblright{}.

\begin{quote}
\pa{ਹਾਦਸੇ} \pa{ਵਿੱਚ} \pa{ਅਮਨ} \pa{ਦੀ} \pa{ਲੱਤ} \pa{ਦੀ} \pa{ਹੱਡੀ} \pa{ਟੁੱਟੀ।}
\end{quote}

\begin{quote}
\emph{hādse vich Aman dī latt dī haḍḍī ṭuṭṭī.}
\end{quote}

\begin{quote}
Aman broke a bone in his leg in the accident.
\end{quote}

A feminine noun.

\begin{cousinweb}
\textbf{\pa{ਹੱਡੀ}} (\emph{haḍḍī}) is inherited from Sanskrit \emph{haḍḍa}, \textquotedblleft{}bone\textquotedblright{}; CDIAL lists Punjabi \emph{haḍḍī}, \textquotedblleft{}bone\textquotedblright{}, under it (CDIAL 13952).
\end{cousinweb}

\begin{grammarlens}[title={What you've built}]
One more word for this chapter.
\end{grammarlens}

\subsection*{Your turn}
\begin{itemize}
\raggedright
  \item \emph{Say it:} \emph{haḍḍī}
  \item \emph{Say it:} \emph{haḍḍī}, once more
  \item \emph{Say it:} \emph{hādse vich Aman dī latt dī haḍḍī ṭuṭṭī.} — Aman broke a bone in his leg in the accident.
  \item \emph{Recall:} say the Punjabi for an ointment, then the Punjabi for the belly, then say \emph{haḍḍī} again
  \item \emph{Write it:} copy \textbf{\pa{ਹੱਡੀ}} once, checking the whole word
\end{itemize}

\begin{tcolorbox}[breakable,colback=teal!4,colframe=teal!35!black,title={Before you move on}]
What does \pa{ਹੱਡੀ} mean? (\textquotedblleft{}A bone\textquotedblright{}.) Say it once more.
\end{tcolorbox}
\section[\texorpdfstring{\pa{ਕਲੇਜਾ} (kalejā)}{ਕਲੇਜਾ (kalejā)}]{\pa{ਕਲੇਜਾ} (kalejā) — the liver}

\label{lesson:PA-C323-kaleja}



\begin{quote}
Before the new one: say the Punjabi for the belly, then the Punjabi for a bone.
\end{quote}

\subsection*{You'll want to know: \pa{ਕਲੇਜਾ}}
\textbf{\pa{ਕਲੇਜਾ}} — \emph{kalejā} — \textquotedblleft{}the liver\textquotedblright{}.

\begin{quote}
\pa{ਕਲੇਜਾ} \pa{ਢਿੱਡ} \pa{ਦੇ} \pa{ਸੱਜੇ} \pa{ਪਾਸੇ} \pa{ਹੁੰਦਾ} \pa{ਹੈ।}
\end{quote}

\begin{quote}
\emph{kalejā ḍhiḍḍ de sajje pāse hundā hai.}
\end{quote}

\begin{quote}
The liver is on the right side of the belly.
\end{quote}

A masculine noun.

\begin{cousinweb}
\textbf{\pa{ਕਲੇਜਾ}} (\emph{kalejā}) is inherited from Sanskrit \emph{kāleyaka}, \textquotedblleft{}a part of the intestines; the liver\textquotedblright{}; CDIAL lists Punjabi \emph{kalejā}, \textquotedblleft{}liver\textquotedblright{}, under it (CDIAL 3103).
\end{cousinweb}

\begin{grammarlens}[title={What you've built}]
One more word for this chapter.
\end{grammarlens}

\subsection*{Your turn}
\begin{itemize}
\raggedright
  \item \emph{Say it:} \emph{kalejā}
  \item \emph{Say it:} \emph{kalejā}, once more
  \item \emph{Say it:} \emph{kalejā ḍhiḍḍ de sajje pāse hundā hai.} — The liver is on the right side of the belly.
  \item \emph{Recall:} say the Punjabi for the belly, then the Punjabi for a bone, then say \emph{kalejā} again
  \item \emph{Write it:} copy \textbf{\pa{ਕਲੇਜਾ}} once, checking the whole word
\end{itemize}

\begin{tcolorbox}[breakable,colback=teal!4,colframe=teal!35!black,title={Before you move on}]
What does \pa{ਕਲੇਜਾ} mean? (\textquotedblleft{}The liver\textquotedblright{}.) Say it once more.
\end{tcolorbox}
\section[\texorpdfstring{\pa{ਪਸਲੀ} (paslī)}{ਪਸਲੀ (paslī)}]{\pa{ਪਸਲੀ} (paslī) — a rib}

\label{lesson:PA-C323-pasli}



\begin{quote}
Before the new one: say the Punjabi for a bone, then the Punjabi for the liver.
\end{quote}

\subsection*{You'll want to know: \pa{ਪਸਲੀ}}
\textbf{\pa{ਪਸਲੀ}} — \emph{paslī} — \textquotedblleft{}a rib\textquotedblright{}.

\begin{quote}
\pa{ਮੇਰੀ} \pa{ਖੱਬੀ} \pa{ਪਸਲੀ} \pa{ਵਿੱਚ} \pa{ਦਰਦ} \pa{ਹੈ।}
\end{quote}

\begin{quote}
\emph{merī khabbī paslī vich darad hai.}
\end{quote}

\begin{quote}
I have a pain in a rib on my left side.
\end{quote}

A feminine noun.

\begin{cousinweb}
\textbf{\pa{ਪਸਲੀ}} (\emph{paslī}) is inherited from Sanskrit \emph{parśu}, \textquotedblleft{}a rib\textquotedblright{}; CDIAL lists Punjabi \emph{paslī}, \textquotedblleft{}rib\textquotedblright{}, under it (CDIAL 7948).
\end{cousinweb}

\begin{grammarlens}[title={What you've built}]
One more word for this chapter.
\end{grammarlens}

\subsection*{Your turn}
\begin{itemize}
\raggedright
  \item \emph{Say it:} \emph{paslī}
  \item \emph{Say it:} \emph{paslī}, once more
  \item \emph{Say it:} \emph{merī khabbī paslī vich darad hai.} — I have a pain in a rib on my left side.
  \item \emph{Recall:} say the Punjabi for a bone, then the Punjabi for the liver, then say \emph{paslī} again
  \item \emph{Write it:} copy \textbf{\pa{ਪਸਲੀ}} once, checking the whole word
\end{itemize}

\begin{tcolorbox}[breakable,colback=teal!4,colframe=teal!35!black,title={Before you move on}]
What does \pa{ਪਸਲੀ} mean? (\textquotedblleft{}A rib\textquotedblright{}.) Say it once more.
\end{tcolorbox}
\section[\texorpdfstring{\pa{ਦੁਖਣਾ} (dukhṇā)}{ਦੁਖਣਾ (dukhṇā)}]{\pa{ਦੁਖਣਾ} (dukhṇā) — to ache}

\label{lesson:PA-C323-dukhna}



\begin{quote}
Before the new one: say the Punjabi for the liver, then the Punjabi for a rib.
\end{quote}

\subsection*{You'll want to know: \pa{ਦੁਖਣਾ}}
\textbf{\pa{ਦੁਖਣਾ}} — \emph{dukhṇā} — \textquotedblleft{}to ache\textquotedblright{}.

\begin{quote}
\pa{ਅੱਜ} \pa{ਮੇਰੀ} \pa{ਪਿੱਠ} \pa{ਬਹੁਤ} \pa{ਦੁਖਦੀ} \pa{ਹੈ।}
\end{quote}

\begin{quote}
\emph{ajj merī piṭṭh bahut dukhdī hai.}
\end{quote}

\begin{quote}
My back aches a lot today.
\end{quote}

Past (it …): \pa{ਉਹ} \pa{ਦੰਦ} \pa{ਰਾਤ} \pa{ਨੂੰ} \pa{ਦੁਖਿਆ}, \emph{oh dand rāt nū̃ dukhiā}.

\begin{cousinweb}
\textbf{\pa{ਦੁਖਣਾ}} (\emph{dukhṇā}) is inherited from Sanskrit \emph{duḥkhati}, \textquotedblleft{}pains\textquotedblright{}; CDIAL lists Punjabi \emph{dukhṇā}, \textquotedblleft{}to hurt\textquotedblright{}, under it (CDIAL 6376).
\end{cousinweb}

\begin{grammarlens}[title={What you've built}]
One more word for this chapter.
\end{grammarlens}

\subsection*{Your turn}
\begin{itemize}
\raggedright
  \item \emph{Say it:} \emph{dukhṇā}
  \item \emph{Say it:} \emph{dukhṇā}, once more
  \item \emph{Say it:} \emph{ajj merī piṭṭh bahut dukhdī hai.} — My back aches a lot today.
  \item \emph{Recall:} say the Punjabi for the liver, then the Punjabi for a rib, then say \emph{dukhṇā} again
  \item \emph{Write it:} copy \textbf{\pa{ਦੁਖਣਾ}} once, checking the whole word
\end{itemize}

\begin{tcolorbox}[breakable,colback=teal!4,colframe=teal!35!black,title={Before you move on}]
What does \pa{ਦੁਖਣਾ} mean? (\textquotedblleft{}To ache\textquotedblright{}.) Say it once more.
\end{tcolorbox}
